\chapter{Preface}

This book has a simple ambition: a single, coherent book of economics
covering everything from the sixth year of school to the end of the
bachelor's degree, written in English for readers anywhere in the world.

The style is concise and rigorous: courses built from definitions, examples,
theorems, methods and diagrams, with careful reasoning whenever it is
accessible at the given level, followed by carefully chosen exercises.
Results stated without full proof are explicitly marked as admitted. Full
solutions to all exercises are collected at the end of the book.

Economics is taught here as a set of models that work, and as the facts that
test them. The course opens with the facts: how living standards, health,
poverty, work, trade and prices have changed over two centuries, read from
charts before any model is drawn. Each later chapter opens with one chart and
explains it. Every model then gets the same four questions: what it assumes,
what it predicts, where it was tested, and where it breaks. The simplest
models --- supply and demand, comparative advantage, the growth model, the
short-run models of the whole economy --- are toy models, and the book says
so; it also says how far each one goes. Where the evidence leaves a question
open, a box sets out the competing estimates, each with its best study and its
weaknesses. Every chart is redrawn from public data by a tested script, and
every number taken from the outside world is traced to its source.

The course comes in five books. The first covers the middle school years,
the second the three years of high school, and the last three the three years
of university. The school books need only the mathematics of the same year of
school; from the first university year, data and econometrics chapters print
Python programs, each taken from a file that has been run and tested, with the
output it actually printed. The two school books never teach the same notion
twice: each idea is taught once, in the year where it first fits, and a later
chapter that needs it opens with a short reminder, \emph{You already know},
before going further. Where a notion belongs to mathematics, the companion
books of the same collection teach it, and this book recalls it before taking
the economist's view. In the school books each year is a part, split into
chapters; each university book is one year, split into parts and chapters.
The earlier the year, the younger the reader it is written for: more
illustrations, more detailed worked steps, and gentler exercises.

\section*{How to read this book}

Statements are numbered within each chapter and color-coded: definitions in
blue, theorems in red, propositions and their relatives in orange, and
methods --- short recipes for standard tasks --- in green. Framed boxes hold
the rest: a model with its four questions; how a number is measured; a dated
episode used as a test of a model; an open question with the estimates on
each side; a reminder of what an earlier chapter taught; and facts that change
from year to year, with the year they are given for. Exercises are graded
from ($\star$) for direct applications to ($\star\star\star$) for harder
problems. A chart's source is named under it, and the full references are
listed in the Data Credits at the front of the book.

\section*{Contributing}

The source of this book lives in a public git repository:
\begin{center}
\url{https://github.com/bvirrion/one-economics-book}
\end{center}
Contributions --- new chapters, corrections, better explanations,
additional exercises --- are welcome; see \texttt{CONTRIBUTING.md} at the
root of the repository.
